\BourbakiExerciseGroup

\begin{enumerate}
\def\labelenumi{\arabic{enumi})}
\tightlist
\item
  A not necessarily commutative ordered monoid M is a (multiplicatively written) monoid equipped with an ordering such that the relation \(x \leqslant y\) implies \(zx \leqslant zy\) and \(xz \leqslant yz\) for all \(z \in M\) . Let G be a not necessarily commutative ordered group. Show that :
\end{enumerate}

\begin{enumerate}
\def\labelenumi{\alph{enumi})}
\item
  If P denotes the set of elements greater than the identity e of G, then P.P = P, \(P \cap P^{-1} = \{e\}\) and \(aPa^{-1} = P\) for all \(a \in G\) . Prove the converse. Find a condition for G to be totally ordered.
\item
  If one of the elements \(\sup (x,y)\) , \(\inf (x,y)\) exists, then so does the other, and
\end{enumerate}

\[
\sup (x, y) = x (\inf (x, y)) ^ {- 1} y = y (\inf (x, y)) ^ {- 1} x.
\]

\begin{enumerate}
\def\labelenumi{\alph{enumi})}
\setcounter{enumi}{2}
\tightlist
\item
  The elements \(\sup(x, e)\) and \(\sup(x^{-1}, e)\) commute. Two coprime elements commute. d) The subgroup \(G'\) generated by the irreducible elements of G is commutative. Deduce that Th. 2 of VI, p.~18 is still valid.
\end{enumerate}

\begin{enumerate}
\def\labelenumi{\arabic{enumi})}
\setcounter{enumi}{1}
\tightlist
\item
  Let E be a lattice with a composition law \((x, y) \mapsto xy\) (not necessarily associative), such that, for all \(a \in E\) , the maps \(x \mapsto ax\) and \(x \mapsto xa\) are order-automorphisms of E. Let \(x_{a}\) (resp. \(_{a}x\) ) denote the element of E defined by \((x_{a})a = x\) (resp. \(a(_{a}x) = x\) ) and suppose that for each \(a \in E\) the maps \(x \mapsto a_{x}\) and \(x \mapsto_{x} a\) are isomorphisms from the ordering of E to the opposite ordering. Show that, for arbitrary x, y, z,
\end{enumerate}

\[
\left(z _ {\inf (x, y)}\right) \cdot x = \left(z _ {y}\right) \cdot \sup (x, y).
\]

\begin{enumerate}
\def\labelenumi{\arabic{enumi})}
\setcounter{enumi}{2}
\item
  Let \(G\) be an ordered group such that the set \(P\) of positive elements is not 0; show that \(G\) is an infinite group, and cannot have a greatest (or a least) element.
\item
  Let G be an ordered group, let P be the set of positive elements of G, and let f be the natural map from G onto a quotient group G/H. For \(f(P)\) to make G/H an ordered group it is necessary and sufficient that \(0 \leqslant y \leqslant x\) and \(x \in H\) imply \(y \in H\) ; then H is called a convex subgroup of G, and G/H is considered as an ordered group in this way. If G is lattice ordered, then a necessary and sufficient condition for G/H to be lattice ordered and for \(f(\sup(x, y)) = \sup(f(x), f(y))\) is that the relations \(|y| \leqslant |x|\) and \(x \in H\) imply \(y \in H\) ; show that this says that H is a filtered convex subgroup. Show that a lattice ordered group with no filtered convex subgroups other than itself and \(\{0\}\) is totally ordered (consider filtered convex subgroups generated by two positive elements of G); * deduce (Gen.~Top., V, p.~25, Ex. 1) that G is then isomorphic to an additive subgroup of the real numbers. *
\item
  Give an example of an ordered group whose ordering is non-discrete, having nonzero elements of finite order (take the quotient of a suitable ordered group by a subgroup H such that \(P \cap H = \{0\}\) ).
\item
  Let G be an ordered group, let P be its set of positive elements. Show that P - P is the largest filtered subgroup of G, and that it is convex. What is the order relation on the quotient?
\item
  In the group \(\mathbf{Z}\) , if \(\mathbf{P}\) is taken to be the set consisting of 0 and the integers \(\geqslant 2\) , then the resulting ordered group is filtered but not lattice ordered (show that the set of elements \(x\) such that \(x \geqslant 0\) and \(x \geqslant 1\) has two distinct minimal elements).
\item
  Let \(x\) be an element of an ordered group such that \(y = \inf(x, 0)\) is defined; if \(n > 0\) is an integer, then \(nx \geqslant 0\) implies \(x \geqslant 0\) (we have
\end{enumerate}

\[
n y = \inf (n x, (n - 1) x, \dots , 0) \geqslant \inf ((n - 1) x, \dots , 0) = (n - 1) y);
\]

and hence \(nx = 0\) implies \(x = 0\) .

\begin{enumerate}
\def\labelenumi{\arabic{enumi})}
\setcounter{enumi}{8}
\item
  In a lattice ordered group G, show that the sum of a family \((H_{\alpha})\) of filtered convex subgroups is a filtered convex subgroup (use VI, p.~12, Cor.).
\item
  Show that, for every finite sequence \((x_{i})\) ( \(1 \leqslant i \leqslant n\) ) of elements of a lattice ordered group \(G\) ,
\end{enumerate}

\[
\sup \left(x _ {i}\right) = \sum_ {i} x _ {i} - \sum_ {i <   j} \inf \left(x _ {i}, x _ {j}\right) + \dots + (- 1) ^ {p + 1} \sum_ {i _ {1} <   i _ {2} <   \dots <   i _ {p}} \inf \left(x _ {i _ {1}}, \dots , x _ {i _ {p}}\right) +   + \dots + (- 1) ^ {n + 1} \inf \left(x _ {1}, \dots , x _ {n}\right).
\]

(Argue by induction based on Prop. 7 (VI, p.~10), using the distributivity of sup over inf.)

\begin{enumerate}
\def\labelenumi{\arabic{enumi})}
\setcounter{enumi}{10}
\item
  Let \((x_{i})\) be a family of n elements of a lattice ordered group G; for each integer k \((0 \leqslant k \leqslant n)\) , let \(d_{k}\) (resp. \(m_{k}\) ) denote the inf (resp. sup) of the \(\binom{n}{k}\) sums of k terms \(x_{i}\) with distinct indices. Show that \(d_{k} + m_{n-k} = x_{1} + x_{2} + \cdots + x_{n}\) .
\item
  Given a lattice ordered group \(\mathbf{G}\) , a subgroup \(\mathbf{H}\) of \(\mathbf{G}\) is called a colattice if for each pair \(x, y\) of elements of \(\mathbf{H}\) , \(\sup_{\mathbf{G}}(x, y) \in \mathbf{H}\) (and so equals \(\sup_{\mathbf{H}}(x, y)\) ).
\end{enumerate}

\begin{enumerate}
\def\labelenumi{\alph{enumi})}
\item
  If \(\mathbf{G} = \mathbf{Q} \times \mathbf{Q} \times \mathbf{Q}\) (where \(\mathbf{Q}\) has the usual ordering), then the subgroup \(\mathbf{H}\) of \((x, y, z)\) such that \(z = x + y\) is a lattice but not a colattice.
\item
  Every filtered convex (Ex. 4) subgroup H of a lattice ordered group G is a colattice.
\item
  Let G be a lattice ordered group and H an arbitrary subgroup of G. Define H' to be the set of infima of finite subsets of H, and H'' to be the set of suprema of finite subsets of H'. Show that H'' is the smallest colattice of G containing H (use the remark on VI, p.~16).
\end{enumerate}

\begin{enumerate}
\def\labelenumi{\arabic{enumi})}
\setcounter{enumi}{12}
\tightlist
\item
  \begin{enumerate}
  \def\labelenumii{\alph{enumii})}
  \tightlist
  \item
    A monoid \(\mathbf{M}\) is said to be lower semi-lattice ordered (or simply semi-lattice ordered) if it is an ordered monoid, if \(\inf(x, y)\) exists for all \(x, y \in \mathbf{M}\) , and if
  \end{enumerate}
\end{enumerate}

\[
\inf (x + z, y + z) = \inf (x, y) + z
\]

for all \(x, y, z \in \mathbf{M}\) . Prove the identities:

\[
\inf (x, z) + \inf (y, z) = \inf (x + y, z + \inf (x, y, z))
\]

\[
\inf (x, y, z) + \inf (x + y, y + z, z + x) = \inf (x, y) + \inf (y, z) + \inf (z, x).
\]

From the first of these relations, deduce that \(x \leqslant z\) and \(y \leqslant z\) imply \(x + y \leqslant z + \inf(x, y)\) .

Show that Prop. 11 and Cor. (VI, pp.~14-15) are valid in a semi-lattice ordered monoid. b) Let M be a semi-lattice ordered monoid and N a submonoid such that \(\inf_{\mathbf{M}}(x,y) = \inf_{\mathbf{N}}(x,y)\) for \(x,y\in \mathbf{N}\) and such that the elements of N are invertible in M. Show that the subgroup G of M consisting of elements \(x - y\) for \(x,y\in \mathbf{N}\) , is lattice ordered.

\begin{enumerate}
\def\labelenumi{\arabic{enumi})}
\setcounter{enumi}{13}
\tightlist
\item
  Show that, in a semi-lattice ordered monoid M,
\end{enumerate}

\[
\inf \left(x _ {i} + y _ {i}\right) \geqslant \inf \left(x _ {i}\right) + \inf \left(y _ {i}\right)
\]

for all finite sequences of n terms \((x_{i})\) , \((y_{i})\) . Deduce the inequalities

\[
(x + y) ^ {+} \leqslant x ^ {+} + y ^ {+}, \quad | x ^ {+} - y ^ {+} | \leqslant | x - y |
\]

in any lattice ordered group.

\begin{enumerate}
\def\labelenumi{\arabic{enumi})}
\setcounter{enumi}{14}
\tightlist
\item
  Show that
\end{enumerate}

\[
\left| x ^ {+} - y ^ {+} \right| + \left| x ^ {-} - y ^ {-} \right| = | x - y |
\]

in a lattice ordered group (notice that \(x - y \leqslant |x - y|\) and \(|x| - |y| \leqslant |x - y|\) ).

\begin{enumerate}
\def\labelenumi{\arabic{enumi})}
\setcounter{enumi}{15}
\tightlist
\item
  Show that
\end{enumerate}

\[
n. \inf (x, y) + \inf (n x, n y) = 2 n. \inf (x, y)
\]

in a semi-lattice ordered monoid (cf.~Ex. 8). Deduce that \(\inf(nx, ny) = n \cdot \inf(x, y)\) if \(\inf(x, y)\) is a regular element.

\begin{enumerate}
\def\labelenumi{\arabic{enumi})}
\setcounter{enumi}{16}
\tightlist
\item
  Let \(x, y, z, t\) be elements of a semi-lattice ordered monoid \(M\) such that \(z \geqslant 0\) and \(t \geqslant 0\) . Prove the inequality
\end{enumerate}

\[
\inf (x + z, y + t) + \inf (x, y) \geqslant \inf (x + z, y) + \inf (x, y + t).
\]

\begin{enumerate}
\def\labelenumi{\arabic{enumi})}
\setcounter{enumi}{17}
\item
  Let \((\mathbf{G}_{\iota})_{\iota \in I}\) be a family of totally ordered groups indexed by a well-ordered set \(I\) ; define an ordering on the direct sum \(\mathbf{G}'\) of the \(\mathbf{G}_{\iota}\) by taking the strictly positive elements to be the nonzero \((x_{\iota})\) such that \(x_{\iota} > 0\) for the first index \(\iota\) such that \(x_{\iota} \neq 0\) . Show that \(\mathbf{G}'\) is totally ordered.
\item
  Let \(G\) be an additive group and \((P_{\alpha})\) a nonempty family of subsets of \(G\) such that \(P_{\alpha} + P_{\alpha} \subset P_{\alpha}\) and \(P_{\alpha} \cap (-P_{\alpha}) = \{0\}\) ; let \(G_{\alpha}\) denote the ordered group obtained from \(G\) by taking \(P_{\alpha}\) as the set of positive elements. Show that \(P + P \subset P\) and \(P \cap (-P) = \{0\}\) , where \(P = \bigcap_{\alpha} P_{\alpha}\) ; if \(H\) is the ordered group obtained from \(G\) by taking \(P\) as the set of
\end{enumerate}

positive elements, show that H is isomorphic to the diagonal subgroup of the product of the \(G_{\alpha}\) .

\begin{enumerate}
\def\labelenumi{\arabic{enumi})}
\setcounter{enumi}{19}
\tightlist
\item
  Let \(G\) be an additive group and let \(P\) be a subset of \(G\) satisfying the following conditions:
\end{enumerate}

\begin{enumerate}
\def\labelenumi{\arabic{enumi}.}
\tightlist
\item
  \(\mathrm{P} + \mathrm{P} = \mathrm{P}\) ; 2. \(\mathrm{P} \cap (-\mathrm{P}) = \{0\}\) ; 3. for every integer \(n > 0\) , \(nx \in \mathrm{P}\) implies \(x \in \mathrm{P}\) (conditions (C)).
\end{enumerate}

\begin{enumerate}
\def\labelenumi{\alph{enumi})}
\item
  If \(a\) is an element of \(G\) such that \(a \notin P\) , show that there exists a subset \(P'\) of \(G\) , satisfying conditions (C), such that \(P \subset P'\) and \(-a \in P'\) (take \(P'\) to be the set of \(x \in G\) such that there exist integers \(m > 0\) and \(n \geqslant 0\) and an element \(y \in P\) with \(mx = -na + y\) .
\item
  Deduce from a) that P is the intersection of those subsets T of G such that \(T + T = T\) , \(T \cap (-T) = \{0\}\) , \(T \cup (-T) = G\) (that is, making G a totally ordered group) and \(P \subset T\) (use Zorn's Lemma).
\item
  In particular if G is an additive group in which every nonzero element has infinite order, then the intersection of all the subsets T of G such that \(T + T = T\) , \(T \cap (-T) = \{0\}\) and \(T \cup (-T) = G\) is \(\{0\}\) .
\end{enumerate}

\begin{enumerate}
\def\labelenumi{\arabic{enumi})}
\setcounter{enumi}{20}
\item
  An ordered group is called lattice orderable if it is isomorphic to a subgroup of a lattice ordered group. Show that it is necessary and sufficient for an ordered group to be lattice orderable that, for every integer \(n > 0\) , the relation \(nx \geqslant 0\) imply \(x \geqslant 0\) (to show that the condition is sufficient, use Ex. 20, b) and 19). Show that every lattice orderable group is isomorphic to a subgroup of a product of totally ordered groups.
\item
  Let G be a lattice orderable group (considered as a Z-module) and E the vector space \(G_{(Q)}\) (II, p.~277); show that there is a unique ordering on E compatible with the additive group structure of E, inducing the given ordering on G, and such that E is lattice orderable.
\item
  Let G be the lattice ordered group \(Z \times Z\) (where Z has the usual ordering) and H the convex subgroup of G generated by \((2, -3)\) ; show that the ordered group G/H is not lattice orderable (cf.~VI, p.~30, Ex. 4).
\item
  Let A be a commutative ring and I the set of all ideals of A. Then \(\mathfrak{a}(\mathfrak{b} + \mathfrak{c}) = \mathfrak{a}\mathfrak{b} + \mathfrak{a}\mathfrak{c}\) for all a, b, c in I, by I, p.~107; in other words I, with the order relation \(a \supset b\) and the law of composition \((\mathfrak{a}, \mathfrak{b}) \mapsto \mathfrak{a}\mathfrak{b}\) , is a semi-lattice ordered monoid (VI, p.~31, Ex. 13), the sup of two elements a, b of I being \(a \cap b\) , and their inf \(a + b\) .
\end{enumerate}

\begin{enumerate}
\def\labelenumi{\alph{enumi})}
\tightlist
\item
  Let K be a field and \(A = K[X, Y]\) the ring of polynomials in two indeterminates over K.
\end{enumerate}

\[
\mathbf {a} = (\mathbf {X}), \mathbf {b} = (\mathbf {Y}), \mathbf {c} = (\mathbf {X} + \mathbf {Y})
\]

\[
\begin{array}{c} (\mathfrak {a} \cap \mathfrak {b}) + \mathfrak {c} \neq (\mathfrak {a} + \mathfrak {c}) \cap (\mathfrak {b} + \mathfrak {c}), \\ (\mathfrak {a} + \mathfrak {b}) \cap \mathfrak {b} \neq (\mathfrak {a} \cap \mathfrak {c}) + (\mathfrak {b} \cap \mathfrak {c}), \\ (\mathfrak {a} \cap \mathfrak {b}) (\mathfrak {a} + \mathfrak {b}) \neq (\mathfrak {a} (\mathfrak {a} + \mathfrak {b})) \cap (\mathfrak {b} (\mathfrak {a} + \mathfrak {b})). \end{array}
\]

\begin{enumerate}
\def\labelenumi{\alph{enumi})}
\setcounter{enumi}{1}
\tightlist
\item
  In the ring A, 1 is a gcd of X and Y, but \((\mathbf{X}) + (\mathbf{Y}) \neq \mathbf{A}\) .
\end{enumerate}

\begin{enumerate}
\def\labelenumi{\arabic{enumi})}
\setcounter{enumi}{24}
\tightlist
\item
  Let I be the semilattice ordered monoid of ideals of a commutative ring A (Ex. 24). For a finite system of congruences \(x \equiv a_i(a_i)\) to admit a solution whenever any two of them admit a common solution (that is to say whenever
\end{enumerate}

\[
a _ {i} \equiv a _ {j} (\mathfrak {a} _ {i} + \mathfrak {a} _ {j})
\]

for every pair of indices \(i, j)\) ; it is necessary and sufficient that each of the laws

\[
(a, b) \mapsto a \cap b, (a, b) \mapsto a + b,
\]

in I be distributive over the other (« Chinese remainder theorem »). One may proceed as follows :

\begin{enumerate}
\def\labelenumi{\alph{enumi})}
\item
  If \((\mathfrak{a}_1 \cap \mathfrak{a}_2) + (\mathfrak{a}_1 \cap \mathfrak{a}_3) = \mathfrak{a}_1 \cap (\mathfrak{a}_2 + \mathfrak{a}_3)\) and if any two of the three congruences \(x \equiv a_i(\mathfrak{a}_i)\) ( \(i = 1, 2, 3\) ) admit a common solution, then the three congruences have a common solution (let \(x_{12}\) be a common solution of \(x \equiv a_1(\mathfrak{a}_1)\) and \(x \equiv a_2(\mathfrak{a}_2)\) and let \(x_{13}\) be a common solution of \(x \equiv a_1(\mathfrak{a}_i)\) and \(x \equiv a_3(\mathfrak{a}_3)\) ; show that the congruences \(x \equiv x_{12}(\mathfrak{a}_1 \cap \mathfrak{a}_2)\) and \(x \equiv x_{13}(\mathfrak{a}_1 \cap \mathfrak{a}_3)\) have a common solution).
\item
  If any system of three congruences, any two of which have a common solution, admits a common solution, show that the following distributive laws hold
\end{enumerate}

\[
\mathfrak {a} + (\mathfrak {b} \cap \mathfrak {c}) = (\mathfrak {a} + \mathfrak {b}) \cap (\mathfrak {a} + \mathfrak {c}), \quad \text { and } \quad \mathfrak {a} \cap (\mathfrak {b} + \mathfrak {c}) = (\mathfrak {a} \cap \mathfrak {b}) + (\mathfrak {a} \cap \mathfrak {c}).
\]

(For the first, note that for all \(x \in (\mathfrak{a} + \mathfrak{b}) \cap (\mathfrak{a} + \mathfrak{c})\) , there exists \(y \in \mathfrak{b} \cap \mathfrak{c}\) such that \(y \equiv x(\mathfrak{a})\) ; for the second, note that for all \(x \in \mathfrak{a} \cap (\mathfrak{b} + \mathfrak{c})\) , there exists \(y \in \mathfrak{a} \cap \mathfrak{b}\) such that \(y \equiv x(\mathfrak{c})\) .) Note that by \(a)\) and \(b)\) , the second distributive law implies the first (cf.~Set Theory, III, p.~217, Ex. 16).

\begin{enumerate}
\def\labelenumi{\alph{enumi})}
\setcounter{enumi}{2}
\item
  Prove the « Chinese remainder theorem » by induction on the number of congruences under consideration, by an analogous argument to that of a).
\item
  * Translation in the case of principal ideal domains. *
\end{enumerate}

\begin{enumerate}
\def\labelenumi{\arabic{enumi})}
\setcounter{enumi}{25}
\item
  Show that the ideal (X) in the monoid of ideals of the polynomial ring K{[}X, Y{]} (where K is a commutative field) satisfies condition (P) of Prop. 14 of VI, p.~17, but is not maximal.
\item
  Let A be the quadratic Z-algebra with basis \((1, e)\) , where \(e^{2} = -5\) . Show that A is an integral domain; in A we have \(9 = 3 \cdot 3 = (2 + e)(2 - e)\) ; show that 3, \((2 + e)\) and \((2 - e)\) are irreducible elements of A, but do not satisfy the condition in Prop. 14 (DIV) of VI, p.~17.
\item
  Let G be a lattice ordered group with positive elements P. Two elements x and y of P are said to be equivalent if any element coprime to one is also coprime to the other; the equivalence classes under this relation are called the threads of P; let \(\bar{x}\) denote the thread containing x.
\end{enumerate}

\begin{enumerate}
\def\labelenumi{\alph{enumi})}
\item
  Let \(\overline{a}\) and \(\overline{b}\) be threads; let \(x, x_1\) be elements of \(\overline{a}\) and \(y, y_1\) be elements of \(\overline{b}\) ; show that if every element coprime to \(x\) is also coprime to \(y\) , then every element coprime to \(x_1\) is also coprime to \(y_1\) . Define a relation in this way between \(\overline{a}\) and \(\overline{b}\) , written \(\overline{a} \geqslant \overline{b}\) ; show that this is an ordering of the set F of threads.
\item
  Show that \(\mathbf{F}\) is a lattice under this ordering, and that
\end{enumerate}

\[
\inf (\bar {a}, \bar {b}) = \overline {{\inf (a , b)}}, \quad \sup (\bar {a}, \bar {b}) = \overline {{\sup (a , b)}} = \overline {{a + b}}
\]

(use Cor. 3 of VI, p.~15). Show that F has a smallest element, namely the thread \(\overline{0} = \{0\}\) .

\begin{enumerate}
\def\labelenumi{\alph{enumi})}
\setcounter{enumi}{2}
\item
  Call two threads \(\overline{a}\) , \(\overline{b}\) coprime if \(\inf (\overline{a},\overline{b}) = \overline{0}\) . Show that if \(\overline{a}\) and \(\overline{b}\) are two threads \(\neq \overline{0}\) , and if \(\overline{a} < \overline{b}\) , then there exists a thread \(\overline{c}\) , coprime to \(\overline{a}\) , such that \(\overline{c} < \overline{b}\) .
\item
  Call a thread \(\bar{m} \neq \bar{0}\) irreducible if it is a minimal element of the set of threads \(\neq \bar{0}\) . Show that an irreducible thread is totally ordered (if \(a\) and \(b\) are two elements of \(\bar{m}\) , consider the threads containing \(a - \inf(a, b)\) and \(b - \inf(a, b)\) and use \(b\) ). Show also that the union of \(\bar{m}, -\bar{m}\) and \(\{0\}\) is a totally ordered subgroup \(\mathrm{H}(\bar{m})\) of \(\mathbf{G}\) .
\item
  Given a family \((\bar{m}_{\iota})\) of distinct irreducible threads of G, show that the subgroup H generated by the union of the \(\bar{m}_{\iota}\) is isomorphic to the direct sum of the groups \(\mathrm{H}(\bar{m}_{\iota})\) (reduce to the case of a finite family, and argue by induction).
\item
  Show that in a product (resp. direct sum) of nonzero totally ordered groups, the elements of irreducible threads are those in which all the coordinates are zero except for one (which is positive). Deduce that an ordered group can be a product (resp. direct sum) of nonzero totally ordered groups in at most one way (in other words the factors are uniquely determined).
\end{enumerate}

\begin{enumerate}
\def\labelenumi{\arabic{enumi})}
\setcounter{enumi}{28}
\tightlist
\item
  An ordered group G is completely lattice ordered if it is filtered and if every nonempty subset of G which is bounded above has a supremum in G. A lower semi-lattice ordered monoid is said to be completely semi-lattice ordered if \(\inf(x_{\alpha})\) exists for every nonempty
\end{enumerate}

family of elements \((x_{\alpha})\) of M which is bounded below, and if

\[
\inf _ {\alpha , \beta} (x _ {\alpha} + y _ {\beta}) = \inf _ {\alpha} (x _ {\alpha}) + \inf _ {\beta} (y _ {\beta})
\]

for any two nonempty families \((x_{\alpha})\) , \((y_{\beta})\) which are bounded below. A completely lattice ordered group is a completely semi-lattice ordered monoid.

\begin{enumerate}
\def\labelenumi{\alph{enumi})}
\item
  Show that a necessary and sufficient condition for a filtered group G to be completely lattice ordered is that every nonempty subset of the set P of positive elements of G admit an infimum in P.
\item
  Every product of completely lattice ordered groups is completely lattice ordered.
\item
  Every filtered convex subgroup of a completely lattice ordered group is completely lattice ordered.
\end{enumerate}

\begin{enumerate}
\def\labelenumi{\arabic{enumi})}
\setcounter{enumi}{29}
\tightlist
\item
  Let G be an ordered group. For every subset A of G let \(m(A)\) (resp. M(A)) denote the set of lower (resp. upper) bounds for A in G; then \(m(A) = -M(-A)\) . a) The relation \(A \subset B\) implies \(m(B) \subset m(A)\) and \(\mathbf{M}(m(A)) \subset \mathbf{M}(m(B))\) .
\end{enumerate}

\begin{enumerate}
\def\labelenumi{\alph{enumi})}
\setcounter{enumi}{1}
\item
  We have \(\mathbf{A} \subset \mathbf{M}(m(\mathbf{A}))\) and \(\mathbf{M}(m(\mathbf{M}(\mathbf{A}))) = \mathbf{M}(\mathbf{A})\) .
\item
  If \((\mathbf{A}_{\iota})\) is an arbitrary family of subsets of \(G\) then \(m\left(\bigcup_{\iota} \mathbf{A}_{\iota}\right) = \bigcap_{\iota} m(\mathbf{A}_{\iota})\) .
\item
  Every set M(B), where B is a nonempty subset bounded above in G, is called a major set in G. For any nonempty subset A bounded below in G, the set M(m(A)) is the smallest major set containing A; denote it \(\langle A\rangle\) ; if \(A\subset B\) then \(\langle A\rangle\subset\langle B\rangle\) ; if A has an infimum a in G, then \(\langle A\rangle=M(a)\) , which we also write \(\langle a\rangle\) .
\item
  If A and B are two nonempty subsets bounded below in G, then \(\langle A + B \rangle = \langle \langle A \rangle + B \rangle\) . Deduce that, in the set \(\mathfrak{M}(G)\) of major sets in G, the map \((A, B) \mapsto \langle A + B \rangle\) is an associative and commutative composition law, for which \(\langle 0 \rangle = P\) is an identity, that the order relation \(A \supset B\) is compatible with this law, and that \(\mathfrak{M}(G)\) is a completely semi-lattice ordered monoid (Ex. 29) under this structure if G is filtered. Moreover, the map \(x \mapsto \langle x \rangle\) from G to \(\mathfrak{M}(G)\) is an isomorphism from the ordered group G onto a subgroup of the monoid \(\mathfrak{M}(G)\) .
\item
  If an element A of \(\mathfrak{M}(G)\) is invertible under the composition law of this monoid, then its inverse is \(\mathbf{M}(-\mathbf{A}) = -m(\mathbf{A})\) (notice that if B is the inverse of A, then the relations \(\mathbf{B} \subset \mathbf{M}(-\mathbf{A})\) and \(\mathbf{A} + \mathbf{M}(-\mathbf{A}) \subset \langle 0 \rangle\) hold, whence \(\langle \mathbf{A} + \mathbf{B} \rangle = \langle \mathbf{A} + \mathbf{M}(-\mathbf{A}) \rangle\) ).
\item
  For an element A of \(\mathfrak{M}(G)\) to be invertible, it is necessary and sufficient that \(x + A \subset A\) imply \(x \geqslant 0\) (write down an expression for \(0 \in \langle A + M(-A) \rangle\) ).
\end{enumerate}

\begin{enumerate}
\def\labelenumi{\arabic{enumi})}
\setcounter{enumi}{30}
\tightlist
\item
  An ordered group G is called archimedean if the only elements \(x \in G\) such that the set of nx (n a positive integer) is bounded below are the positive elements of G.
\end{enumerate}

\begin{enumerate}
\def\labelenumi{\alph{enumi})}
\item
  The monoid \(\mathfrak{M}(G)\) of major sets in an ordered group G is a group if and only if G is archimedean (use Ex. 30, g) and d). If in addition G is filtered, then \(\mathfrak{M}(G)\) is a completely lattice ordered group.
\item
  Deduce that a filtered group G is isomorphic to a subgroup of a completely lattice ordered group if and only if G is archimedean.
\end{enumerate}

\begin{enumerate}
\def\labelenumi{\arabic{enumi})}
\setcounter{enumi}{31}
\tightlist
\item
  \begin{enumerate}
  \def\labelenumii{\alph{enumii})}
  \tightlist
  \item
    Let G be a completely lattice ordered group, and H a subgroup of G. For every major subset A of the group H, let \(x_{\mathbf{A}}\) denote the infimum of A in G; show that the map \(\mathbf{A} \mapsto x_{\mathbf{A}}\) from \(\mathfrak{M}(\mathbf{H})\) to G is injective (show that A is the set of \(y \in \mathbf{H}\) such that \(y \geqslant x_{\mathbf{A}}\) ).
  \end{enumerate}
\end{enumerate}

\begin{enumerate}
\def\labelenumi{\alph{enumi})}
\setcounter{enumi}{1}
\item
  For every subset B of H, bounded below in H, let \(\langle B\rangle\) denote the major set generated by B in H (an element of \(\mathfrak{M}(H)\) ). If for any subset B of H, bounded below in H, the equation \(\inf B = \inf \langle B \rangle\) holds (infima taken in G), show that the map \(A \mapsto x_{A}\) is an isomorphism of the ordered group \(\mathfrak{M}(H)\) onto a subgroup of G (cf.~Ex. 30, e)).
\item
  If every element of G is the infimum of some subset of H, show that for any subset B of H, bounded below in H, we have inf B = inf \(\langle B\rangle\) (infima taken in G). (Notice that this holds if inf B ∈ H; in general, consider an element \(x \in G\) such that \(x + inf B \in H\) and use the fact that x is the infimum of some subset of G, as well as Ex. 30, e).
\end{enumerate}

\(d) ^{*}\) Let \(G\) be the completely lattice ordered group \(\mathbf{Z} \times \mathbf{Z} \times \mathbf{R}\) . Let \(\theta > 0\) be an irrational number; let \(H\) be the subgroup of \((x, y, z)\) such that \(\theta(z - x) + y = 0\) . Show that there exists no isomorphism of the ordered group \(\mathfrak{M}(H)\) onto any subgroup of \(G\) which restricts to the identity map on \(H\) (notice that \(\mathfrak{M}(H)\) is isomorphic to \(K = Z \times R\) ; consider the subgroup of \(K\) consisting of elements \(u\) such that, for every integer \(n > 0\) , there exists \(v \in K\) with \(nv = u\) , and the analogous subgroup of \(G\) ). *

\begin{enumerate}
\def\labelenumi{\arabic{enumi})}
\setcounter{enumi}{32}
\tightlist
\item
  \begin{enumerate}
  \def\labelenumii{\alph{enumii})}
  \tightlist
  \item
    A totally ordered group G is archimedean if and only if, for any pair \(x > 0\) , \(y > 0\) of elements of G, there exists an integer \(n > 0\) such that \(y \leqslant nx\) .
  \end{enumerate}
\end{enumerate}

* b) Any totally ordered, archimedean, completely lattice ordered group G is isomorphic to \(\{0\}\) , Z or R (ruling out the first two cases, fix a \textgreater{} 0 in G, and make each \(x \in G\) correspond to the infimum of the rational numbers p/q such that \(qx \leq pa\) ).

\begin{enumerate}
\def\labelenumi{\alph{enumi})}
\setcounter{enumi}{2}
\item
  Deduce that every archimedean totally ordered group G is isomorphic to a subgroup of R (notice that \(\mathfrak{M}(G)\) is totally ordered). *
\item
  The lexicographic product \(Z \times Z\) is totally ordered and not archimedean.
\end{enumerate}

\begin{enumerate}
\def\labelenumi{\arabic{enumi})}
\setcounter{enumi}{33}
\item
  Let \(G = Z^{N}\) be the product of a countably infinite family of totally ordered groups Z, and let \(H = Z^{(N)}\) be the filtered convex subgroup, the direct sum of the factors of G. Show that the lattice ordered group G/H (VI, p.~30, Ex. 4) is not archimedean, although G and H are completely lattice ordered.
\item
  A major set A in an ordered group G is said to be finitely generated if there exists a finite set F such that \(A = \langle F \rangle\) . We say that G is semi-archimedean if every finitely generated major set is invertible in the monoid \(\mathfrak{M}(G)\) . Every lattice ordered (resp. archimedean) group is semi-archimedean (Ex. 31, a)).
\end{enumerate}

\begin{enumerate}
\def\labelenumi{\alph{enumi})}
\item
  Show that every filtered semi-archimedean group is lattice ordered (if \(nx \geqslant 0\) , consider the major set \(\langle F \rangle\) , where \(F = \{0, x, \ldots, (n - 1)x\}\) , and use Ex. 30, g) and the fact that it is invertible).
\item
  * Let K be the lexicographic product \(R \times R\) , and G the (usual) product group \(K \times R\) . Let \(\theta\) be an irrational number such that \(0 < \theta < 1\) , and let H be the subgroup generated by \(((1, 0), 0)\) , \(((\theta, 0), \theta)\) and \(((0, x), 0)\) where x runs through R. Show that H, a subgroup of a product of two totally ordered groups, is not semi-archimedean (consider the major set generated by the two elements \(((1,0),0)\) and \(((\theta ,0),\theta)\) of H, and show that it is not invertible). *
\item
  Let G be a semi-archimedean filtered group; show that the submonoid of \(\mathfrak{M}(G)\) generated by the finitely generated major sets and their inverses is a lattice ordered group (cf.~VI, p.~31, Ex. 13, b)).
\item
  * Let \(\theta > 0\) be an irrational number, and let \(G\) be the group \(\mathbf{Z} \times \mathbf{Z}\) in which the set \(P\) of positive elements is taken to be the set of \((x, y)\) such that \(x \geqslant 0\) and \(y \geqslant \theta x\) . Show that \(G\) is an archimedean group, but that the inverse in \(\mathfrak{M}(G)\) of a finitely generated major set, not of the form \(\langle a \rangle\) , is not finitely generated. *
\end{enumerate}

\begin{enumerate}
\def\labelenumi{\arabic{enumi})}
\setcounter{enumi}{35}
\tightlist
\item
  Let \(\mathbf{G}\) be a filtered group with positive elements \(\mathbf{P}\) . Then \(\mathbf{G}\) is isomorphic to a group of the form \(\mathbf{Z}^{(I)}\) if and only if there exists a map \(x \mapsto d(x)\) from \(\mathbf{P}\) into \(\mathbf{N}\) such that, for all \(b\) , either \(b \geqslant a\) or there exists an element \(c\) in the major set \(\langle a, b \rangle\) such that \(d(c) < d(a)\) . (Show that this condition is satisfied in \(\mathbf{Z}^{(I)}\) ) by the map \(d(x_{i}) = \sum_{i} x_{i}\) (when \(x_{i} \geqslant 0\) for all \(i\) ).
\end{enumerate}

Conversely, if the condition is satisfied, show that every major set \(\mathbf{A} \subset \mathbf{P}\) has the form \(\langle a \rangle\) , by taking an element \(a\) of \(\mathbf{A}\) such that \(d(a) \leqslant d(x)\) for all \(x \in \mathbf{A}\) . Apply Th. 2 of VI, p.~18.)
% semantic-fragments: legacy-input-seam
\relax{}
